\section{Discussion}\label{sec:discussion}

\paragraph{A familiar bias, not yet a persistent one.} The simplest reading of the general-election results is that
election markets show the favourite--longshot bias long documented in racetrack and sports betting
\citep{griffith1949,thaler1988,snowberg2010}: favourites are priced too cheaply and long shots too dearly. We cannot yet
establish it as a persistent feature of election markets. The evidence rests on two general-election nights, 2024 and
2026, and two nights that point the same way may still be a coincidence; allowing for states moving together, the
direct test of persistence is not significant ($p = 0.13$ on Polymarket). US general elections come only every two
years, so settling the question needs a much longer record. Primaries also do not fit the pattern. Kalshi's primaries
were strongly overconfident, the opposite of the betting-market bias, and they trade far less than general elections
(Section~\ref{sec:data}). Whatever drives the bias, it seems to depend on the kind of race and who trades it, not only
on the market being a betting market.

\paragraph{Information and hedging rather than profit.} Even if the bias persists, little is left after costs: buying
favourites above 70\% earned a few percent per bet on the midterms, and the more profitable rules were found by
searching the same nights on which they were judged (Section~\ref{sec:backtest}). The greater value of these markets
probably lies elsewhere. Their prices are fast, public forecasts that have rivalled polls
\citep{berg2008,rothschild2009}, and the results here hint that a small correction learned from past general elections can make
them slightly more accurate. Their contracts also let investors hedge exposure to an event directly, such as a holding in a solar-energy
fund before the 2024 presidential election; such hedges look practical for small positions, though thin liquidity
limits them at institutional scale \citep{kirillov2026hedge}.

\paragraph{Rules shape who trades.} Regulation and politics decide which venue attracts which traders, which contracts
can be listed at all, and which markets draw volume. Kalshi could list US election contracts only after winning its
case against the CFTC in 2024; Polymarket blocked US users for most of the period studied and re-entered the US market
in 2025 (Section~\ref{sec:data}). The calibration measured here belongs to that particular period of fast growth and
changing rules. A new US administration, or a change in how regulators treat election contracts, could alter who trades
on each platform, which races are listed and how heavily each is traded, and with them the patterns reported in this
paper.
